\scanpage{28}
On a
\[
(15)\quad\left\{\begin{gathered}
M_0\text{ et }\mathbb{F}_2^{\underline{\omega}}
\text{ sont orthogonaux l'un de l'autre}\\
\text{dans }\uline{\operatorname{End}}(M_0)\text{ pour la forme }\operatorname{Tr}(uv)
\end{gathered}\right.
\]
\[
(16)\quad\deleted{\unsure{\mathbb{F}_2}}\ 
\underset{\operatorname{Ker}(\operatorname{End}(M_0)\xrightarrow{\mathrm{tr}}\mathbb{F}_2)}
{\mathfrak{sl}(M_0)}
=M_0\oplus\mathbb{F}_2\operatorname{id}_{M_0}
\]
\[
(17)\quad\mathbb{F}_2^{\underline{\omega}}\cap
\mathfrak{sl}(M_0)
\simeq\mathbb{F}_2\cdot\operatorname{id}_{M_0}\quad\ldots\qquad]
\]
Ainsi, la donnée d'un torseur $P_M$ sous \inserted{$V_J(M_0)\simeq$}$\operatorname{End}(M_0)$\deleted{\unsure{$a$}}, équivaut à celle d'un couple de torseurs, \deleted{sur} l'un sous $M_0$, l'autre sous $\mathbb{F}_2^{\underline{\omega}}$, soient $I_M$ et $S_M$.

\uline{Théorème} Le foncteur $M\mapsto(M_0,I_M,S_M)$
\diagram{\begin{tikzcd}
(18)\quad\substack{\mathbb{Z}/4\mathbb{Z}\text{-modules}\\\text{libres de rang }2}
\arrow[r] & \text{triplets }(M_0,I_{\deleted{M}},S_{\deleted{M}})
\end{tikzcd}}
\deleted{où $M_0$ un vect. de rg $2$}
avec
\[
\left\{\begin{array}{l}
M_0\text{ vect. de rg }2\text{ sur }\mathbb{F}_2\\
I\quad M_0\text{-torseur}\\
S\quad\text{torseur sous }\mathbb{F}_2^{\underline{\omega}(M_0)}
\end{array}\right.
\]
est une équivalence de groupoïdes.

\margin{Équivalent essentiellement à l'isomorphisme
\[
\operatorname{GL}(2,\mathbb{Z}/4\mathbb{Z})
\simeq\operatorname{GL}(2,\mathbb{F}_2)\cdot M_2(\mathbb{F}_2)
\]
(produit $\frac12$ direct), plus l'isomorphisme de $\unsure{\mathcal{G}_3|\mathfrak{S}_3}$-modules ($\unsure{\mathcal{G}_3|\mathfrak{S}_3}\simeq\operatorname{GL}(2,\mathbb{F}_2)$)
\[
M_2(\mathbb{F}_2)\simeq\mathbb{F}_2^2\times\mathbb{F}_2^{\underline{\omega}_0},
\]
où $\operatorname{GL}(2,\mathbb{F}_2)$ opère sur $\mathbb{F}_2^2$ tautologiquement, et sur $\mathbb{F}_2^{\underline{\omega}_0}$ via op. sur $\underline{\omega}_0$ par signature\ldots On a l'inclusion
\[
\operatorname{SL}(2,\mathbb{Z}/4\mathbb{Z})
\unsure{\subset|=}\operatorname{GL}(2,\mathbb{F}_2)\cdot
\deleted{\unsure{M_2(\mathbb{F}_2)}}
\]
\deleted{\unsure{$\operatorname{GL}(2,\mathbb{F}_2)\ltimes\mathbb{F}_2^2\times\mathbb{F}_2$}}
\[
\{g\cdot u\in\unsure{\mathcal{G}_3|\mathfrak{S}_3}\cdot M_2(\mathbb{F}_2)
\mid\operatorname{sg}(g)\underline{e}(\operatorname{Tr}u)=1\}
\]
où $\underline{e}:\deleted{\unsure{\mathbb{F}_2}}\mathbb{F}_2
\xrightarrow{\sim}\{\pm1\}$.}

\uline{Corollaire} On a un isom. de suites exactes canoniques ($\forall M$ libre de rg $2$ sur $\Lambda=\mathbb{Z}/4\mathbb{Z}$)
\diagram{\begin{tikzcd}
1\arrow[r,dash]&\operatorname{End}_{\mathbb{F}_2}(M_0)\arrow[r]&
\operatorname{GL}(M)\arrow[r]&
\operatorname{GL}(M_0)\simeq\unsure{\mathcal{G}_J|\mathfrak{S}_J}\arrow[r,dash]&1\\
(19)\quad1\arrow[r,dash]&
M_0\times\mathbb{F}_2^{\underline{\omega}}\arrow[u,"\sim"]\arrow[r]&
\operatorname{Aut}(M_0,I_M,S_M)\arrow[u,"\sim"]\arrow[r]&
\operatorname{GL}(M_0)\arrow[u,equal]\arrow[r,dash]&1.
\end{tikzcd}}
Ici l'hom $\operatorname{End}(M_0)\longrightarrow\operatorname{GL}(M)$ est donné par la condition que le composé avec l'hom $\operatorname{End}_\Lambda(M)\longrightarrow\operatorname{End}_{\mathbb{F}_2}(M_0)$ soit donné par
\[
u\longmapsto\operatorname{id}_M+2u.
\]
% archive foliation: 27
